% Target: Advances in Applied Clifford Algebras (Springer)
% Format: Standard Springer svjour3 / or article class
% ~17 pages, Clifford-algebra emphasis
% NO overlap with Paper 1 (CICM 2026, methodology/taxonomy/Dollard case study)
%
% This paper is MATHEMATICAL CONTENT: the cross-domain identities themselves.
% Target: AACA (Springer). Using article class for submission.
% Springer reformats to svjour3 during production if accepted.
\documentclass[11pt]{article}
\usepackage[margin=1in]{geometry}

\usepackage{amsmath,amssymb,amsthm}
\usepackage{mathtools}
\usepackage{booktabs}
\usepackage{listings}
\usepackage{xcolor}
\usepackage{hyperref}
\usepackage{cleveref}
\usepackage{enumitem}
\usepackage{tikz-cd}

% Theorem environments
\newtheorem{theorem}{Theorem}[section]
\newtheorem{proposition}[theorem]{Proposition}
\newtheorem{identity}[theorem]{Identity}
\newtheorem{analogy}[theorem]{Analogy}
\newtheorem{corollary}[theorem]{Corollary}
\newtheorem{remark}{Remark}[section]

% Compatibility definitions for Springer svjour3 commands used in body
\newenvironment{acknowledgements}{\section*{Acknowledgements}}{}
\providecommand{\journalname}[1]{}
\providecommand{\titlerunning}[1]{}
\providecommand{\authorrunning}[1]{}
\providecommand{\institute}[1]{}
\providecommand{\email}[1]{\texttt{#1}}
\providecommand{\keywords}[1]{%
  \begingroup\providecommand{\and}{}%
  \renewcommand{\and}{$\cdot$}%
  \par\medskip\noindent\textbf{Keywords:} #1\endgroup}

% Lean 4 listing style
\lstdefinelanguage{Lean4}{
  keywords={theorem, lemma, def, noncomputable, import, by, have, exact, rw,
            simp, ring, norm_num, apply, intro, induction, unfold, show,
            refine, suffices, convert, push_cast, field_simp, ext, linarith,
            nlinarith, native_decide, omega,
            private, where, with, match, if, then, else, let, in, do, return,
            structure, class, instance, open, section, namespace, end, variable,
            axiom, sorry, example, calc, fun, set_option, maxHeartbeats},
  morekeywords=[2]{Cl30, Mat3, SL3, isPureBivector, mkBivector, comm, smul, zero,
            grade0, grade1, grade2, grade3},
  sensitive=true,
  morecomment=[l]{--},
  morecomment=[s]{/-}{-/},
  morestring=[b]",
}

\lstset{
  language=Lean4,
  basicstyle=\ttfamily\scriptsize,
  keywordstyle=\color{blue}\bfseries,
  keywordstyle=[2]\color{purple},
  commentstyle=\color{gray}\itshape,
  stringstyle=\color{red},
  breaklines=true,
  columns=flexible,
  frame=single,
  xleftmargin=1em,
  numbers=left,
  numberstyle=\tiny\color{gray},
  numbersep=3pt,
  aboveskip=0.5em,
  belowskip=0.5em,
  literate=
    {forall}{{$\forall$}}1
    {exists}{{$\exists$}}1,
}

% Notation shortcuts
\newcommand{\so}{\mathfrak{so}}
\newcommand{\su}{\mathfrak{su}}
\newcommand{\g}{\mathfrak{g}}
\newcommand{\h}{\mathfrak{h}}
\newcommand{\Cl}{\mathrm{Cl}}
\newcommand{\ad}{\mathrm{ad}}
\newcommand{\Tr}{\mathrm{Tr}}
\newcommand{\spec}{\mathrm{spec}}
\newcommand{\ZZ}{\mathbb{Z}}
\newcommand{\RR}{\mathbb{R}}
\newcommand{\CC}{\mathbb{C}}
\newcommand{\NN}{\mathbb{N}}

\journalname{Advances in Applied Clifford Algebras}

\begin{document}

\title{Character Decomposition from Fortescue to Cartan--Weyl: \\
  Unifying Symmetrical Components and Root Space Decomposition \\
  via Clifford Algebras}

\titlerunning{Character Decomposition from Fortescue to Cartan--Weyl}

\authorrunning{I.\,M.\ Arnoldy}

\author{Ian M. Arnoldy\thanks{Independent Researcher.
  E-mail: \texttt{ian@arnoldy.me}}}

\institute{
  Independent Researcher \\
  \email{ian@arnoldy.me}
}

\date{March 10, 2026}

\maketitle

\begin{abstract}
The grade-2 elements of a Clifford algebra $\Cl(n,0)$ form the
Lie algebra $\so(n)$ under the commutator bracket. This classical
observation has algebraic consequences that connect three seemingly
disparate decomposition methods: Fortescue's symmetrical components
(electrical engineering, 1918), the Cartan--Weyl root space
decomposition (Lie theory, 1913--1925), and the principal grading by
Coxeter elements \cite{coxeter1934}. We prove that the first two are instances of
character decomposition on compact groups at different levels of
generality. The Fortescue matrix is the character table of the cyclic
group $\ZZ_N$; the Cartan--Weyl decomposition is a Fortescue-type
decomposition of the adjoint representation; and the Coxeter grading
decomposes a simple Lie algebra into $h$ eigenspaces ($h$ = Coxeter
number) via a generalized DFT. We further prove that Clifford grade
selection rules force Hamiltonians built from $\so(n)$ generators into
block-tridiagonal form, establishing a structural constraint on the
spectrum of such operators. All concrete instances---including $\so(3)$
structure constants, Casimir eigenvalues, Killing form negativity, grade
projections in $\Cl(3,0)$, and the complete Cartan--Weyl decomposition
of $\su(3)$---are machine-verified in Lean~4 with zero \texttt{sorry}
gaps. We delineate two additional structural analogies (the mass gap as a
spectral stop band, and the Coxeter grading as generalized Fortescue)
that share the same mathematical parent but require qualification before
being called provable identities.

\keywords{Clifford algebras \and symmetrical components \and
Cartan--Weyl decomposition \and character theory \and Fortescue \and
grade selection rules \and Lean~4}
\end{abstract}

% ============================================================================
\section{Introduction}
\label{sec:intro}
% ============================================================================

In 1918, building on Steinmetz's phasor methods
\cite{steinmetz1893}, Charles L.\ Fortescue presented a method for
decomposing an unbalanced $N$-phase electrical system into $N$ balanced ``symmetrical
components'' \cite{fortescue1918}. The method, which reduces coupled
polyphase problems to decoupled single-phase problems, became a
cornerstone of power systems engineering. Its mathematical content is
the discrete Fourier transform (DFT) on the cyclic group $\ZZ_N$:
the Fortescue matrix $A_N$ with entries $A_{ps} = \omega_N^{ps}$
(where $\omega_N = e^{2\pi i/N}$) is precisely the character table
of $\ZZ_N$.

Meanwhile, between 1913 and 1925, \'Elie Cartan and Hermann Weyl
developed the root space decomposition of semisimple Lie algebras
\cite{cartan1913,weyl1925}. For a semisimple Lie algebra $\g$ with
Cartan subalgebra $\h$, the adjoint representation decomposes into
eigenspaces under $\h$:
\begin{equation}\label{eq:cartan-weyl}
  \g = \h \oplus \bigoplus_{\alpha \in \Phi} \g_\alpha,
\end{equation}
where $\Phi$ is the root system and $\g_\alpha = \{X \in \g : [H, X] =
\alpha(H) X \text{ for all } H \in \h\}$.

These two decomposition methods were developed in complete isolation.
Fortescue's work belongs to electrical engineering; Cartan--Weyl theory
belongs to pure mathematics. Despite a century of development in both
fields, we are not aware of any published work that systematically
identifies the precise mathematical relationship between them.

This paper establishes that relationship. Our main results are:

\begin{enumerate}[leftmargin=*]
\item \textbf{Identity~1} (\Cref{sec:peter-weyl}): The Fortescue
  decomposition is a special case of the Peter--Weyl theorem.
  Specifically, Fortescue's symmetrical components are the character
  decomposition of the regular representation of $\ZZ_N$, and the
  Peter--Weyl theorem generalizes this from $\ZZ_N$ to all compact
  groups.

\item \textbf{Identity~2} (\Cref{sec:cartan-weyl-fortescue}): The
  Cartan--Weyl root space decomposition \eqref{eq:cartan-weyl} is a
  Fortescue-type decomposition of the adjoint representation. The
  Cartan subalgebra plays the role of the ``zero sequence,'' positive
  roots play the role of ``positive sequences,'' and negative roots
  play the role of ``negative sequences.'' We verify this concretely
  for $\su(2)$ (from $\Cl(3,0)$) and $\su(3)$.

\item \textbf{Grade selection rules} (\Cref{sec:grade-selection}):
  In $\Cl(n,0)$, the product of two grade-2 elements lies in grade
  $0 \oplus 2$ (no odd-grade components). This forces any Hamiltonian
  built from $\so(n)$ generators into a block-tridiagonal structure in
  the grade basis, constraining its spectral properties via
  continued-fraction expansions of the resolvent.

\item \textbf{Two analogies} (\Cref{sec:analogy-stopband,sec:analogy-coxeter}):
  We identify two further structural parallels---the mass gap as a
  spectral stop band, and the Coxeter principal grading as a
  generalized Fortescue decomposition---that share the same
  mathematical parent (character theory / spectral decomposition) but
  fall short of provable identities. We explain precisely where each
  analogy breaks.
\end{enumerate}

All concrete computations---structure constants, Casimir eigenvalues,
Killing form negativity, grade projections, and the complete
Cartan--Weyl decomposition of $\su(3)$---are machine-verified in
Lean~4 \cite{demoura2021lean4} using mathlib \cite{mathlib}, with zero
\texttt{sorry} gaps. The proof files are available from the
corresponding author upon reasonable request \cite{uft-github}.

\paragraph{What this paper is not.}
We make no claims about the Yang--Mills mass gap, quantum field theory,
or the physics of gauge theories. Our contribution is purely algebraic:
identifying the common mathematical structure underlying decomposition
methods from different fields, and proving the relevant algebraic facts
in a theorem prover. We are careful throughout to distinguish
\emph{identities} (same mathematics, different notation) from
\emph{analogies} (parallel structure, different scope).

% ============================================================================
\section{Preliminaries}
\label{sec:prelim}
% ============================================================================

\subsection{Fortescue's Symmetrical Components}

Let $\omega_N = e^{2\pi i/N}$ be a primitive $N$th root of unity. The
\emph{Fortescue matrix} is the $N \times N$ matrix
\begin{equation}\label{eq:fortescue-matrix}
  A_N = \frac{1}{\sqrt{N}} \begin{pmatrix}
    1 & 1 & 1 & \cdots & 1 \\
    1 & \omega_N & \omega_N^2 & \cdots & \omega_N^{N-1} \\
    1 & \omega_N^2 & \omega_N^4 & \cdots & \omega_N^{2(N-1)} \\
    \vdots & & & \ddots & \vdots \\
    1 & \omega_N^{N-1} & \omega_N^{2(N-1)} & \cdots & \omega_N^{(N-1)^2}
  \end{pmatrix}.
\end{equation}
This is, up to normalization, the DFT matrix. It is unitary:
$A_N^* A_N = I_N$.

Given an $N$-phase signal $\mathbf{x} = (x_0, x_1, \ldots, x_{N-1})^T$,
the \emph{symmetrical components} are
\begin{equation}\label{eq:symm-comp}
  \mathbf{s} = A_N^* \, \mathbf{x},
\end{equation}
and the inverse transform recovers the phase quantities:
$\mathbf{x} = A_N \, \mathbf{s}$.

The key engineering property: if the system impedance matrix $Z$ is
\emph{circulant} (i.e., commutes with cyclic permutation), then $A_N$
diagonalizes it:
$A_N^* Z A_N = \mathrm{diag}(Z_0, Z_1, \ldots, Z_{N-1})$.
The $Z_k$ are the \emph{sequence impedances}: the coupled $N$-phase
problem decouples into $N$ independent single-phase problems.

For $N = 3$ (three-phase power), the sequences are traditionally called
``zero'' ($k = 0$), ``positive'' ($k = 1$), and ``negative'' ($k = 2$).

\subsection{Character Theory of Finite Abelian Groups}

A \emph{character} of a finite group $G$ is a homomorphism
$\chi : G \to \CC^*$. The characters of the cyclic group $\ZZ_N$ are
$\chi_k(n) = \omega_N^{kn}$ for $k = 0, 1, \ldots, N-1$. The
\emph{character table} of $\ZZ_N$ is the matrix
$T_{kn} = \chi_k(n) = \omega_N^{kn}$---which is exactly the
(unnormalized) Fortescue matrix \eqref{eq:fortescue-matrix}.

This is the central observation: the Fortescue decomposition \emph{is}
the character decomposition of $\ZZ_N$.

\subsection{Clifford Algebras and $\so(n)$}

The Clifford algebra $\Cl(n,0)$ is generated by $n$ elements
$e_1, \ldots, e_n$ satisfying $e_i e_j + e_j e_i = 2\delta_{ij}$
\cite{atiyah-bott-shapiro1964,hestenes1966}.
It has dimension $2^n$ and is $\ZZ_2$-graded into even and odd parts.
More finely, it admits a direct-sum decomposition by \emph{grade}:
grade-$k$ elements are spanned by products of $k$ distinct generators.

\begin{proposition}[Classical; see {\cite[Ch.~3]{lounesto2001}}]
\label{prop:grade2-lie}
The grade-2 elements of\/ $\Cl(n,0)$ form a Lie algebra under the
commutator bracket $[A, B] = AB - BA$. This Lie algebra is isomorphic
to $\so(n)$.
\end{proposition}

The isomorphism sends the bivector $e_{ij} = e_i e_j$ ($i < j$) to the
generator of rotation in the $ij$-plane
\cite{gallier2020}. The dimension is
$\binom{n}{2} = n(n-1)/2$, matching $\dim \so(n)$.

\subsection{The Cartan--Weyl Decomposition}

For a semisimple Lie algebra $\g$ of rank $r$
\cite{humphreys1972,hall2015}, the
\emph{Cartan subalgebra} $\h$ is a maximal abelian subalgebra of
dimension $r$. The \emph{roots} are the nonzero linear functionals
$\alpha \in \h^*$ for which the root space
$\g_\alpha = \{X \in \g : [H, X] = \alpha(H) X, \, \forall H \in \h\}$
is nonzero. The decomposition \eqref{eq:cartan-weyl} splits $\g$ into
$\h$ and the root spaces, with
\begin{equation}\label{eq:dim-decomp}
  \dim \g = \dim \h + |\Phi| = r + |\Phi|.
\end{equation}

The roots come in pairs $\pm \alpha$, giving a further decomposition
into positive and negative roots:
$\g = \mathfrak{n}^- \oplus \h \oplus \mathfrak{n}^+$.

% ============================================================================
\section{Identity~1: Peter--Weyl Generalizes Fortescue}
\label{sec:peter-weyl}
% ============================================================================

\subsection{Statement}

\begin{identity}[Fortescue = Character Decomposition of $\ZZ_N$]
\label{id:fortescue-peter-weyl}
Fortescue's symmetrical component decomposition
\eqref{eq:symm-comp} is the decomposition of a signal
$\mathbf{x} \in \CC^N$ into irreducible representations of the
cyclic group $\ZZ_N$. Equivalently, the Fortescue matrix $A_N$ is the
character table of\/ $\ZZ_N$.
\end{identity}

\subsection{Proof}

The cyclic group $\ZZ_N$ has exactly $N$ irreducible representations,
all one-dimensional, given by the characters
$\chi_k : \ZZ_N \to \CC^*$, $\chi_k(n) = \omega_N^{kn}$. The space
$\CC^N$, viewed as the regular representation of $\ZZ_N$ (where the
generator acts by cyclic permutation), decomposes as
\begin{equation}\label{eq:zn-decomp}
  \CC^N = \bigoplus_{k=0}^{N-1} V_k,
  \qquad V_k = \{v \in \CC^N : \sigma(v) = \omega_N^k v\},
\end{equation}
where $\sigma$ is the cyclic shift operator. The projection onto $V_k$
is
\begin{equation}\label{eq:projection}
  P_k = \frac{1}{N} \sum_{n=0}^{N-1} \overline{\chi_k(n)} \, \sigma^n
      = \frac{1}{N} \sum_{n=0}^{N-1} \omega_N^{-kn} \, \sigma^n.
\end{equation}
The matrix implementing all projections simultaneously is
$T_{kn} = \omega_N^{kn}$, which is (up to normalization) the Fortescue
matrix \eqref{eq:fortescue-matrix}. Thus the Fortescue decomposition
\emph{is} the decomposition \eqref{eq:zn-decomp}. \qed

\subsection{Generalization via Peter--Weyl}

The Peter--Weyl theorem \cite{peter-weyl1927} states that for any
compact group $G$, the space $L^2(G)$ decomposes as
\begin{equation}\label{eq:peter-weyl}
  L^2(G) \cong \widehat{\bigoplus_{\pi \in \hat{G}}}
    \dim(\pi) \cdot V_\pi,
\end{equation}
where the sum runs over equivalence classes of irreducible unitary
representations. For $G = \ZZ_N$, all irreps are one-dimensional
($\dim(\pi) = 1$), $\hat{G} = \ZZ_N$, and
\eqref{eq:peter-weyl} reduces to \eqref{eq:zn-decomp}. Thus:

\begin{proposition}\label{prop:peter-weyl-fortescue}
Fortescue's symmetrical component decomposition is the Peter--Weyl
theorem specialized to $G = \ZZ_N$. The Peter--Weyl theorem extends
Fortescue's method from cyclic groups to all compact Lie groups.
\end{proposition}

This is not merely a restatement. It has a concrete consequence:
any function on a compact Lie group $G$ (not just $\ZZ_N$) admits a
``symmetrical component decomposition'' into irreducible
representations. The Fourier coefficients generalize the sequence
components, and the diagonalization of circulant matrices generalizes
to the diagonalization of convolution operators on $G$.

\subsection{Machine Verification}

The roots-of-unity foundation is verified in Lean~4:

\begin{lstlisting}[caption={$\omega_N^N = 1$: the central theorem of Fortescue's decomposition (excerpt from \texttt{polyphase\_formula.lean}).}]
theorem omega_pow_N (N : Nat)
    (hN : (N : Complex) != 0) :
    omega N ^ N = 1 := by
  simp only [omega]
  rw [exp_nat_mul]
  have h_exp : (N : Complex) *
    (2 * Real.pi * Complex.I / (N : Complex))
    = 2 * Real.pi * Complex.I :=
    omega_exponent_simplifies N hN
  rw [h_exp]
  exact Complex.exp_two_pi_mul_I
\end{lstlisting}

Additional verified results include: the balanced-phase property
$1 + \omega_3 + \omega_3^2 = 0$, the four-phase specialization
$\omega_4 = i$ (connecting to $\ZZ_4$), periodicity
$\omega_N^{N+k} = \omega_N^k$, and the group law
$\omega_N^a \cdot \omega_N^b = \omega_N^{a+b}$.

% ============================================================================
\section{Identity~2: Cartan--Weyl as Sequence Decomposition}
\label{sec:cartan-weyl-fortescue}
% ============================================================================

\subsection{Statement}

\begin{identity}[Cartan--Weyl = Fortescue on the Adjoint Representation]
\label{id:cartan-weyl-fortescue}
The Cartan--Weyl root space decomposition of a semisimple Lie algebra
$\g$ is the eigendecomposition of the adjoint action of the Cartan
subalgebra $\h$. This is a Fortescue-type decomposition in the
following precise sense:
\begin{itemize}[leftmargin=*]
\item The Cartan subalgebra $\h$ (eigenvalue $0$ under $\ad_H$ for all
  $H \in \h$) plays the role of the \emph{zero sequence}.
\item The positive root spaces $\g_\alpha$ ($\alpha > 0$) play the role
  of \emph{positive sequences}.
\item The negative root spaces $\g_{-\alpha}$ play the role of
  \emph{negative sequences}.
\end{itemize}
\end{identity}

\subsection{The Case $\su(2)$ from $\Cl(3,0)$}

The bivectors of $\Cl(3,0)$ are $\{e_{12}, e_{13}, e_{23}\}$, forming
$\so(3) \cong \su(2)$ under the commutator. Define the angular momentum
generators $L_k = \frac{1}{2} e_{jl}$ (with appropriate index
convention). The Cartan--Weyl basis is:
\begin{align}
  H &= L_3 = \tfrac{1}{2} e_{12}
    &\text{(Cartan, ``zero sequence'')}, \nonumber \\
  E_+ &= L_1 + i L_2
    &\text{(positive root, ``positive sequence'')}, \label{eq:su2-cw} \\
  E_- &= L_1 - i L_2
    &\text{(negative root, ``negative sequence'')}. \nonumber
\end{align}

The adjoint eigenvalues are:
\begin{equation}\label{eq:su2-eigenvalues}
  [H, E_+] = +E_+, \qquad [H, E_-] = -E_-, \qquad [H, H] = 0.
\end{equation}
This is a decomposition into three ``sequences'' with
eigenvalues $\{+1, -1, 0\}$ under the adjoint action of the Cartan
element---exactly the Fortescue pattern for $N = 3$ reduced to its
essential structure: one zero, one positive, one negative.

The underlying Clifford algebra structure is verified in Lean~4:

\begin{lstlisting}[caption={$\so(3)$ structure constants from $\Cl(3,0)$ (excerpt from \texttt{grade2\_lie\_algebra.lean}).}]
-- [e12, e13] = -2*e23
theorem comm_e12_e13 :
    comm e12 e13 = smul (-2) e23 := by
  ext <;> simp [comm, e12, e13, e23,
    mul, add, neg, smul] <;> norm_num

-- [e12, e23] = 2*e13
theorem comm_e12_e23 :
    comm e12 e23 = smul 2 e13 := by
  ext <;> simp [comm, e12, e23, e13,
    mul, add, neg, smul] <;> norm_num

-- Jacobi identity: [A,[B,C]] + cyc = 0
theorem jacobi (A B C : Cl30) :
    comm A (comm B C) + comm B (comm C A)
    + comm C (comm A B) = zero := by
  ext <;> simp [comm, mul_def, add_def,
    neg_def, mul, add, neg, zero] <;> ring
\end{lstlisting}

\subsection{The Case $\su(3)$: Complete Verification}

The Lie algebra $\su(3)$ has dimension $8 = 3^2 - 1$, rank $2$, and
$6$ roots. In the Chevalley basis $\{H_1, H_2, E_1, F_1, E_2, F_2,
E_3, F_3\}$ \cite{georgi1999,fulton1991}, the Cartan--Weyl
decomposition is:
\begin{equation}\label{eq:su3-decomp}
  \su(3) = \underbrace{\langle H_1, H_2 \rangle}_{\text{Cartan (rank 2)}}
    \oplus \underbrace{\langle E_1, E_2, E_3 \rangle}_{\text{positive roots}}
    \oplus \underbrace{\langle F_1, F_2, F_3 \rangle}_{\text{negative roots}}.
\end{equation}
The dimension count is $8 = 2 + 3 + 3$.

The adjoint eigenvalues form the root system of type $A_2$. We list
the complete eigenvalue table, which is machine-verified:

\begin{table}[htbp]
\centering
\caption{Adjoint eigenvalues for $\su(3)$. Each entry is the eigenvalue
  $\alpha(H_i)$ in $[H_i, E_\alpha] = \alpha(H_i) E_\alpha$.}
\label{tab:su3-roots}
\begin{tabular}{@{}lccl@{}}
\toprule
Generator & $\ad_{H_1}$ & $\ad_{H_2}$ & Root \\
\midrule
$H_1$ & $0$ & $0$ & zero \\
$H_2$ & $0$ & $0$ & zero \\
$E_1$ & $+2$ & $-1$ & $\alpha_1$ \\
$F_1$ & $-2$ & $+1$ & $-\alpha_1$ \\
$E_2$ & $-1$ & $+2$ & $\alpha_2$ \\
$F_2$ & $+1$ & $-2$ & $-\alpha_2$ \\
$E_3$ & $+1$ & $+1$ & $\alpha_1 + \alpha_2$ \\
$F_3$ & $-1$ & $-1$ & $-(\alpha_1 + \alpha_2)$ \\
\bottomrule
\end{tabular}
\end{table}

The Fortescue--Cartan--Weyl identity is now concrete:
\begin{itemize}[leftmargin=*]
\item \textbf{Zero sequence} ($\alpha = 0$): $\{H_1, H_2\}$, the
  simultaneously diagonalizable generators. In power systems, the zero
  sequence represents in-phase components.
\item \textbf{Positive sequences} ($\alpha > 0$): $\{E_1, E_2, E_3\}$,
  the raising operators. These carry positive ``frequency'' under the
  adjoint action.
\item \textbf{Negative sequences} ($\alpha < 0$): $\{F_1, F_2, F_3\}$,
  the lowering operators.
\end{itemize}

The Cartan matrix of type $A_2$,
\begin{equation}\label{eq:cartan-matrix}
  A = \begin{pmatrix} 2 & -1 \\ -1 & 2 \end{pmatrix},
  \qquad \det A = 3 \neq 0,
\end{equation}
encodes the root geometry. It is the analog of the Fortescue matrix
restricted to simple roots. The nonzero determinant (verified:
\texttt{cartan\_matrix\_determinant}) proves the root system is
non-degenerate, just as the invertibility of the Fortescue matrix
guarantees the decomposition is unique.

Additional verified structure:
\begin{itemize}[leftmargin=*]
\item $\mathfrak{n}^+ = \langle E_1, E_2, E_3 \rangle$ is a closed
  subalgebra: $[E_1, E_2] = E_3$ and all other brackets among
  $E_i$ vanish (\texttt{positive\_root\_closed}).
\item Similarly $\mathfrak{n}^-$ is closed (\texttt{negative\_root\_closed}).
\item Each pair $(E_\alpha, F_\alpha, [E_\alpha, F_\alpha])$ forms an
  $\su(2)$ triple: $[E_1, F_1] = H_1$, $[E_2, F_2] = H_2$,
  $[E_3, F_3] = H_1 + H_2$.
\item All $\binom{8}{2} = 28$ independent brackets are verified
  (\texttt{bracket\_count}).
\end{itemize}

\begin{lstlisting}[caption={Root eigenvalue verification and subalgebra closure for $\su(3)$ (excerpts from \texttt{su3\_cartan\_weyl.lean}).}]
-- [H1, E1] = 2*E1 (root alpha_1(H1) = 2)
theorem root_alpha1_H1 :
    comm H1 E1 = smul 2 E1 := h1e1

-- [H2, E1] = -E1 (root alpha_1(H2) = -1)
theorem root_alpha1_H2 :
    comm H2 E1 = neg E1 := by
  ext <;> simp [comm, H2, E1, neg]

-- Positive root subalgebra is CLOSED
theorem positive_root_closed (A B : SL3)
    (hA : A.h1 = 0 /\ A.h2 = 0 /\
          A.f1 = 0 /\ A.f2 = 0 /\ A.f3 = 0)
    (hB : B.h1 = 0 /\ B.h2 = 0 /\
          B.f1 = 0 /\ B.f2 = 0 /\ B.f3 = 0) :
    (comm A B).h1 = 0 /\ (comm A B).h2 = 0 /\
    (comm A B).f1 = 0 /\ (comm A B).f2 = 0 /\
    (comm A B).f3 = 0 := by ...
\end{lstlisting}

\subsection{The General Pattern}

The parallel between Fortescue and Cartan--Weyl extends to any
semisimple Lie algebra. Let $\g$ have rank $r$ and $|\Phi^+|$ positive
roots. Then:
\begin{equation}\label{eq:general-decomp}
  \dim \g = r + 2|\Phi^+|
    = \underbrace{r}_{\text{zero}} + \underbrace{|\Phi^+|}_{\text{positive}}
    + \underbrace{|\Phi^+|}_{\text{negative}}.
\end{equation}

\begin{table}[htbp]
\centering
\caption{Dimension decomposition for classical Lie algebras compared to
  Fortescue.}
\label{tab:decomp-comparison}
\begin{tabular}{@{}lcccl@{}}
\toprule
Algebra & Zero & Pos. & Neg. & Total \\
\midrule
Fortescue ($\ZZ_3$) & 1 & 1 & 1 & 3 \\
$\su(2)$ & 1 & 1 & 1 & 3 \\
$\su(3)$ & 2 & 3 & 3 & 8 \\
$\su(5)$ & 4 & 10 & 10 & 24 \\
$\so(10)$ & 5 & 20 & 20 & 45 \\
$\so(14)$ & 7 & 42 & 42 & 91 \\
\bottomrule
\end{tabular}
\end{table}

The common mathematical parent is harmonic analysis on compact groups.
Fortescue uses characters of $\ZZ_N \subset U(1)$; Cartan--Weyl uses
characters of the maximal torus $T^r \subset G$. Both decompose a
representation into eigenspaces of a maximal abelian subgroup.

% ============================================================================
\section{The Casimir Operator and Killing Form}
\label{sec:casimir}
% ============================================================================

The quadratic Casimir operator $C_2 = \sum_a T_a T_a$ (summing over an
orthonormal basis of $\g$ with respect to the Killing form) commutes
with all elements of $\g$ and acts as a scalar on each irreducible
representation (Schur's lemma). The Casimir eigenvalue provides an
algebraic measure of ``size'' for each representation.

\subsection{The $\so(3)$ Casimir from $\Cl(3,0)$}

In $\Cl(3,0)$, the Casimir of the bivector algebra is
\begin{equation}\label{eq:casimir-cl30}
  C_2 = e_{12}^2 + e_{13}^2 + e_{23}^2.
\end{equation}
Since each $e_{ij}^2 = -1$ in $\Cl(3,0)$ (from $e_i^2 = +1$ and
anticommutativity), we get $C_2 = -3$ (as a scalar in the algebra).

In the fundamental (3-dimensional) matrix representation, the Casimir
acts as
\begin{equation}\label{eq:casimir-fund}
  C_2 = L_1^2 + L_2^2 + L_3^2 = -2 I_3,
\end{equation}
where $L_i$ are the standard $\so(3)$ generators.

Both facts are machine-verified:

\begin{lstlisting}[caption={Casimir computation in $\Cl(3,0)$ and the $3 \times 3$ representation (excerpts from \texttt{grade2\_lie\_algebra.lean} and \texttt{casimir\_eigenvalues.lean}).}]
-- In Cl(3,0): C2 = -3 (scalar)
theorem bivector_sum_of_squares_val :
    bivector_sum_of_squares =
    smul (-3) (1 : Cl30) := by
  ext <;> simp [bivector_sum_of_squares,
    e12, e13, e23, mul, add, one, smul]
    <;> norm_num

-- In Mat3: C2 = -2*I3
theorem casimir_fund_val :
    casimir_fund = smul (-2) id3 := by
  ext <;> simp [casimir_fund, L1, L2, L3,
    mul, add, smul, id3] <;> norm_num
\end{lstlisting}

\subsection{The Killing Form and Compactness}

The Killing form $B(X, Y) = \Tr(\ad_X \circ \ad_Y)$ restricted to
$\so(3)$ is negative definite:

\begin{theorem}[Machine-verified]\label{thm:killing-negative}
For $H = aL_1 + bL_2 + cL_3 \in \so(3)$ with $(a, b, c) \neq (0,0,0)$:
\begin{equation}\label{eq:killing}
  \Tr(H^2) = -2(a^2 + b^2 + c^2) < 0.
\end{equation}
\end{theorem}

\noindent The proof in Lean~4 proceeds by direct computation of the
trace, followed by the algebraic fact that $a^2 + b^2 + c^2 > 0$
when $(a,b,c) \neq 0$.

\begin{lstlisting}[caption={Killing form negativity (from \texttt{casimir\_spectral\_gap.lean}).}]
theorem killing_form_negative (a b c : Real)
    (h : !(a = 0 /\ b = 0 /\ c = 0)) :
    tr (so3_element_squared a b c) < 0 := by
  rw [so3_sq_trace]
  have : a^2 + b^2 + c^2 > 0 := by
    by_contra hle; push_neg at hle
    have ha : a = 0 := by nlinarith
      [sq_nonneg a, sq_nonneg b, sq_nonneg c]
    have hb : b = 0 := by nlinarith ...
    have hc : c = 0 := by nlinarith ...
    exact h (And.intro ha (And.intro hb hc))
  linarith
\end{lstlisting}

The negative definiteness of the Killing form is the algebraic signature
of \emph{compactness}. A compact semisimple Lie group has all
finite-dimensional representations unitarizable, discrete Casimir
spectrum, and---crucially for applications---the Casimir eigenvalue
provides a lower bound on representation ``energies.''

\subsection{Casimir Eigenvalues for $\so(n)$}

For the fundamental representation of $\so(n)$ with the standard
normalization $\Tr(T_a T_b) = -\delta_{ab}$ (antisymmetric matrices),
the Casimir eigenvalue is $C_2(\text{fund}) = -(n-1)$. With the
half-normalization $\Tr(T_a T_b) = -\frac{1}{2}\delta_{ab}$, this
becomes $c_2(\text{fund}) = (n-1)/2$.

\begin{table}[htbp]
\centering
\caption{Casimir eigenvalues for $\so(n)$ in three representations.}
\label{tab:casimir}
\begin{tabular}{@{}lccc@{}}
\toprule
$n$ & $c_2(\text{fund})$ & $c_2(\text{adj})$ & $c_2(\text{spinor})$ \\
\midrule
10 & $9/2$ & $16$ & $45/4$ \\
14 & $13/2$ & $24$ & $91/4$ \\
\bottomrule
\end{tabular}
\end{table}

The minimum nonzero Casimir eigenvalue (the fundamental) provides the
tightest algebraic constraint on representation sizes. For $\so(14)$:
$c_2(\text{fund}) = (n-1)/2 = 13/2$, verified in Lean~4.

The Casimir centrality---$[C_2, T_a] = 0$ for all generators
$T_a$---is verified for all three $\so(3)$ generators, confirming
that $C_2$ acts as a scalar on each irreducible component (Schur's
lemma).

% ============================================================================
\section{Analogy~1: Mass Gap as Stop Band}
\label{sec:analogy-stopband}
% ============================================================================

We now turn to two structural parallels that share the character-theoretic
parent of Identities~1 and~2 but require honest qualification.

\subsection{The Spectral Correspondence}

In quantum mechanics, the Hamiltonian $\hat{H}$ generates time evolution
$U(t) = e^{-i\hat{H}t}$. The \emph{resolvent} (Green's function) is
\begin{equation}\label{eq:resolvent}
  G(z) = (z - \hat{H})^{-1},
\end{equation}
whose poles are exactly $\spec(\hat{H})$.

In circuit theory, the impedance $Z(s)$ is a Laplace-domain transfer
function whose poles correspond to natural frequencies. The
\emph{mass gap} condition
\begin{equation}\label{eq:mass-gap}
  \spec(\hat{H}) = \{0\} \cup [\Delta, \infty), \qquad \Delta > 0,
\end{equation}
states that $G(z)$ has no poles in $(0, \Delta)$. In engineering
language, this is a \emph{stop band}: no resonances between the vacuum
state ($z = 0$) and the first excited state ($z = \Delta$).

\begin{analogy}[Mass Gap = Stop Band]\label{an:stopband}
The operator-theoretic content of the mass gap condition
\eqref{eq:mass-gap} is identical to the statement that a transfer
function has a stop band from $0$ to $\Delta$. The resolvent
\eqref{eq:resolvent} is mathematically a transfer function.
\end{analogy}

\subsection{Where the Analogy Breaks}

The spectral formulation is exact: the resolvent \emph{is} a transfer
function, and the mass gap \emph{is} a stop band, in the precise sense
of spectral theory of self-adjoint operators. However, three critical
distinctions prevent this from being a \emph{proof tool}:

\begin{enumerate}[leftmargin=*]
\item \textbf{Finite vs.\ infinite dimensions.} Circuit transfer
  functions are rational (finitely many poles). The Yang--Mills
  Hamiltonian acts on an infinite-dimensional Hilbert space with
  potentially continuous spectrum.

\item \textbf{Linearity.} Sequence networks in power systems work
  because Maxwell's equations are linear. The Yang--Mills Hamiltonian
  includes the non-abelian self-interaction term $[A_\mu, A_\nu]$,
  which breaks superposition.

\item \textbf{Existence.} Writing down \eqref{eq:resolvent} presupposes
  that $\hat{H}$ exists as a self-adjoint operator on a suitable Hilbert
  space. For interacting quantum Yang--Mills theory in four dimensions,
  this existence is itself unproven---it is part of the Clay Millennium
  Prize problem \cite{jaffe-witten2000}.
\end{enumerate}

We rate this as a genuine structural analogy (the mathematical
formalism is identical at the spectral-theoretic level) that
nevertheless cannot substitute for the analytical work of proving
\eqref{eq:mass-gap} for a specific interacting quantum field theory.

% ============================================================================
\section{Analogy~2: Fortescue--Coxeter Principal Grading}
\label{sec:analogy-coxeter}
% ============================================================================

\subsection{The Coxeter Element and Principal Grading}

Let $\g$ be a simple Lie algebra of rank $r$ with Weyl group $W$. A
\emph{Coxeter element} $\varphi \in W$ is a product of all simple
reflections (one from each simple root), taken in any order. All
such products are conjugate and have the same order $h$, the
\emph{Coxeter number} \cite{humphreys1990}.

The Coxeter element $\varphi$ acts on $\g$ via the adjoint
representation, and its eigenvalues are $\{\omega_h^{m_j} : j = 1,
\ldots, r\}$, where $m_1, \ldots, m_r$ are the \emph{exponents} of
$\g$ and $\omega_h = e^{2\pi i/h}$. The principal grading decomposes
$\g$ into $h$ eigenspaces:
\begin{equation}\label{eq:coxeter-grading}
  \g = \bigoplus_{k=0}^{h-1} \g_k,
  \qquad \g_k = \{X \in \g : \varphi(X) = \omega_h^k X\}.
\end{equation}

\begin{analogy}[Coxeter Grading $\approx$ Fortescue with Period $h$]
\label{an:coxeter}
The principal grading \eqref{eq:coxeter-grading} is a DFT
decomposition with period $h$. For $\su(N)$ (type $A_{N-1}$), the
Coxeter number is $h = N$, so the principal grading \emph{is}
an $N$-phase Fortescue decomposition with the correct period.
\end{analogy}

\subsection{Where the Periods Match}

For the $A$-series:
\begin{equation}\label{eq:A-series}
  \su(N): \quad h = N, \quad
  \text{exponents} = \{1, 2, \ldots, N-1\}.
\end{equation}
The Coxeter decomposition \eqref{eq:coxeter-grading} has exactly $N$
eigenspaces, labeled by $\omega_N^k$ for $k = 0, \ldots, N-1$. This
is the $N$-phase Fortescue decomposition on the Lie algebra,
not on signal space.

\subsection{Where the Periods Diverge}

For other Dynkin types, $h \neq \dim(\text{fundamental representation})$,
and the Fortescue period does not directly match the representation
dimension. For example:

\begin{table}[htbp]
\centering
\caption{Coxeter numbers and exponents for selected algebras.}
\label{tab:coxeter}
\begin{tabular}{@{}llcl@{}}
\toprule
Type & Algebra & $h$ & Exponents \\
\midrule
$A_2$ & $\su(3)$ & 3 & $\{1, 2\}$ \\
$A_4$ & $\su(5)$ & 5 & $\{1, 2, 3, 4\}$ \\
$D_5$ & $\so(10)$ & 8 & $\{1, 3, 4, 5, 7\}$ \\
$D_7$ & $\so(14)$ & 12 & $\{1, 3, 5, 6, 7, 9, 11\}$ \\
\bottomrule
\end{tabular}
\end{table}

For $\so(14) = D_7$: the Coxeter number is $h = 12$ (not 14), and the
exponents $\{1, 3, 5, 6, 7, 9, 11\}$ are not consecutive. The
Fortescue interpretation requires a 12-phase system with non-uniform
mode occupancy. The grading \eqref{eq:coxeter-grading} is still a
DFT decomposition, but with period 12, not 14. The rank-7 Cartan
subalgebra occupies $\g_0$, and the 84 root vectors are distributed
among $\g_1$ through $\g_{11}$ according to the exponent structure,
with the dimension check $7 + 84 = 91 = \dim \so(14)$.

\subsection{Why This Is an Analogy, Not an Identity}

The Coxeter grading \eqref{eq:coxeter-grading} is a character
decomposition under $\ZZ_h \subset W$ (the cyclic group generated by
the Coxeter element), so it \emph{is} a DFT decomposition. However:

\begin{enumerate}[leftmargin=*]
\item The Coxeter element acts on the Lie algebra via the adjoint
  representation, not on a signal space. The engineering content of
  Fortescue (decoupling coupled equations via impedance diagonalization)
  does not directly apply.

\item For non-$A$ types, the exponents are not consecutive, so
  not all $\ZZ_h$ eigenspaces are occupied. The decomposition has
  ``missing harmonics.''

\item The Coxeter number $h$ enters Lie theory through the formula
  $|W| = \prod_{i=1}^r (m_i + 1)$ and the dual Coxeter number
  $\check{h}$ controls the level of the affine Lie algebra. These
  connections go beyond the Fourier-analytic content of Fortescue.
\end{enumerate}

We rate this as a precise analogy: same abstract structure
(character decomposition under a cyclic group), same mathematical
parent (representation theory), but different scope and not
directly useful for the engineering applications that motivate
Fortescue's method.

% ============================================================================
\section{Grade Selection Rules and Block-Tridiagonal Structure}
\label{sec:grade-selection}
% ============================================================================

This section contains the most novel algebraic result of the paper:
the observation that Clifford grade selection rules impose a structural
constraint on operators built from Lie algebra generators.

\subsection{Grade Projections in $\Cl(n,0)$}

The algebra $\Cl(n,0)$ decomposes into $n+1$ grade components:
\begin{equation}\label{eq:grade-decomp}
  \Cl(n,0) = \bigoplus_{k=0}^{n} \Cl^k(n,0),
  \qquad \dim \Cl^k(n,0) = \binom{n}{k}.
\end{equation}
Each grade-$k$ subspace is spanned by products of $k$ distinct
generators. For $\Cl(3,0)$: grades $0, 1, 2, 3$ have dimensions
$1, 3, 3, 1$ (total $8 = 2^3$).

The grade projections $\pi_k : \Cl(n,0) \to \Cl^k(n,0)$ are verified
to be idempotent ($\pi_k^2 = \pi_k$), mutually orthogonal
($\pi_j \circ \pi_k = 0$ for $j \neq k$), and exhaustive
($\sum_k \pi_k = \mathrm{id}$).

\subsection{The Selection Rule}

\begin{theorem}[Grade Selection Rule; Machine-verified]
\label{thm:selection-rule}
Let $A, B \in \Cl^2(3,0)$ be pure bivectors. Then their Clifford
product $AB$ has no grade-1 or grade-3 component:
\begin{equation}\label{eq:selection-rule}
  \pi_1(AB) = 0, \qquad \pi_3(AB) = 0.
\end{equation}
Equivalently, $AB \in \Cl^0(3,0) \oplus \Cl^2(3,0)$.
\end{theorem}

\begin{proof}
Direct computation in the 8-component representation of $\Cl(3,0)$.
The Lean~4 proof expands the product using the multiplication table,
substitutes the vanishing components (a bivector has zero scalar,
vector, and pseudoscalar parts), and verifies that the resulting vector
and pseudoscalar components are identically zero by the \texttt{ring}
tactic.
\end{proof}

\begin{lstlisting}[caption={Grade selection rule (from \texttt{block\_tridiagonal.lean}).}]
-- (grade 2) * (grade 2) has no odd grades
theorem grade2_times_grade2_selection
    (A B : Cl30)
    (hA : isPureBivector A)
    (hB : isPureBivector B) :
    (A * B).v1 = 0 /\ (A * B).v2 = 0 /\
    (A * B).v3 = 0 /\ (A * B).ps = 0 := by
  obtain (hAs, hAv1, hAv2, hAv3, hAps) := hA
  obtain (hBs, hBv1, hBv2, hBv3, hBps) := hB
  refine (And.intro ?_ (And.intro ?_
    (And.intro ?_ ?_))) <;>
  simp [mul_def, mul, hAs, hAv1, hAv2,
    hAv3, hAps, hBs, hBv1, hBv2,
    hBv3, hBps] <;> ring
\end{lstlisting}

\subsection{The General Grade Shifting Rule}

More generally, the Clifford product of a grade-$p$ element with a
grade-$q$ element has components only in grades
$|p - q|, |p - q| + 2, \ldots, p + q$ (subject to the constraint
$\leq n$). For $p = q = 2$: the product lies in grades $0$, $2$,
and (when $n \geq 4$) grade~$4$.

The commutator $[A, B] = AB - BA$ of two bivectors is therefore a
pure bivector (grade 2). The grade-0 part of the Clifford product
is symmetric: $\langle AB \rangle_0 = \langle BA \rangle_0$
(it equals the scalar product $A \cdot B$). The grade-4 part is
likewise symmetric: $\langle AB \rangle_4 = A \wedge B = B \wedge A
= \langle BA \rangle_4$ (since the exterior product of two
2-forms is symmetric). Hence both cancel in the antisymmetric
combination $AB - BA$, leaving only grade~2, which confirms
\Cref{prop:grade2-lie}.

The complete grade-shifting behavior of bivectors is verified in Lean~4:
\begin{itemize}[leftmargin=*]
\item $e_{12} \cdot 1 = e_{12}$ \quad (grade $0 \to 2$),
\item $e_{12} \cdot e_1 = -e_2$ \quad (grade $1 \to 1$, rotation),
\item $e_{12} \cdot e_3 = e_{123}$ \quad (grade $1 \to 3$, shift by $+2$),
\item $e_{12} \cdot e_{23} = e_{13}$ \quad (grade $2 \to 2$, no shift),
\item $e_{12} \cdot e_{123} = -e_3$ \quad (grade $3 \to 1$, shift by $-2$).
\end{itemize}

\subsection{Block-Tridiagonal Hamiltonian Structure}

Consider an operator $\hat{H}$ built from grade-2 elements
(i.e., from $\so(n)$ generators). In the Clifford grade basis, the
matrix elements $\langle \psi_p | \hat{H} | \psi_q \rangle$ are
nonzero only when $|p - q| \in \{0, 2\}$. This means $\hat{H}$
has \emph{block-tridiagonal} structure:

\begin{equation}\label{eq:block-tridiag}
  \hat{H} \sim \begin{pmatrix}
    H_{00} & H_{02} & 0 & 0 & \cdots \\
    H_{20} & H_{22} & H_{24} & 0 & \cdots \\
    0 & H_{42} & H_{44} & H_{46} & \cdots \\
    0 & 0 & H_{64} & H_{66} & \ddots \\
    \vdots & & & \ddots & \ddots
  \end{pmatrix},
\end{equation}
where $H_{pq}$ is the block connecting grade-$p$ and grade-$q$ sectors.

\begin{theorem}[Block-Tridiagonal Structure]\label{thm:block-tridiag}
Any operator on $\Cl(n,0)$ that is a polynomial in grade-2 elements
has, in the grade basis, a banded matrix with bandwidth $2$. The
even-grade and odd-grade sectors decouple:
\begin{equation}\label{eq:decouple}
  \Cl(n,0) = \Cl^{\mathrm{even}}(n,0) \oplus \Cl^{\mathrm{odd}}(n,0),
\end{equation}
and $\hat{H}$ preserves each summand.
\end{theorem}

\begin{proof}
The even subalgebra $\Cl^+(n,0) \coloneqq \Cl^{\mathrm{even}}(n,0)
= \bigoplus_{k \text{ even}} \Cl^k$ is a subalgebra (closed under
multiplication), and any product of grade-2 elements lies in
$\Cl^+$. Hence $\hat{H}$ maps even grades
to even grades and odd grades to odd grades, giving the
block-diagonal decomposition \eqref{eq:decouple}. Within each parity
sector, the bandwidth-2 banding follows from the grade selection rule
(\Cref{thm:selection-rule}).
\end{proof}

\subsection{Spectral Consequences}

Block-tridiagonal matrices have a well-developed spectral theory via
\emph{continued fractions}. The resolvent of a block-tridiagonal
operator admits the expansion
\begin{equation}\label{eq:continued-fraction}
  G(z) = \left(z - a_0 - \frac{b_1^2}{z - a_1 -
    \frac{b_2^2}{z - a_2 - \cdots}}\right)^{-1},
\end{equation}
where $a_k$ and $b_k$ are derived from the diagonal and off-diagonal
blocks. The poles of \eqref{eq:continued-fraction} are the eigenvalues
of $\hat{H}$, and the structure of the continued fraction constrains
their distribution.

For $\Cl(14,0)$ with $\so(14)$ generators: the algebra has dimension
$2^{14} = 16384$, with 15 grades of dimensions
$\binom{14}{k}$. The even subalgebra $\Cl^+(14,0) \cong \Cl(13,0)$
has dimension $2^{13} = 8192$. Furthermore, $\Cl(14,0) \cong M(128,
\RR)$ (the algebra of $128 \times 128$ real matrices), providing $128$
primitive idempotents that further decompose the representation space.

The block-tridiagonal structure, combined with the $128$-sector
idempotent decomposition, constrains the sparsity pattern of any
Hamiltonian built from $\so(14)$ generators. This structural
information goes beyond what the Casimir eigenvalues alone encode.

% ============================================================================
\section{Machine Verification}
\label{sec:verification}
% ============================================================================

All concrete algebraic results in this paper are machine-verified in
Lean~4 \cite{demoura2021lean4} with mathlib \cite{mathlib}, compiled
with zero \texttt{sorry} gaps. \Cref{tab:proofs} summarizes the
relevant proof files.

\begin{table*}[t]
\centering
\caption{Machine-verified theorems relevant to this paper. All files
  compile with zero \texttt{sorry} gaps.}
\label{tab:proofs}
\small
\begin{tabular}{@{}llp{6cm}@{}}
\toprule
\textbf{File} & \textbf{Key Theorem} & \textbf{Content} \\
\midrule
\texttt{polyphase\_formula.lean}
  & \texttt{omega\_pow\_N}
  & $\omega_N^N = 1$; roots of unity, DFT kernel \\
  & \texttt{cube\_roots\_sum\_zero}
  & $1 + \omega_3 + \omega_3^2 = 0$; balanced-phase property \\
  & \texttt{omega4\_eq\_I}
  & $\omega_4 = i$; connecting $\ZZ_4$ to complex plane \\
\midrule
\texttt{grade2\_lie\_algebra.lean}
  & \texttt{comm\_pure\_bivector}
  & $[\text{grade 2}, \text{grade 2}] \subseteq \text{grade 2}$; Lie closure \\
  & \texttt{jacobi}
  & $[A,[B,C]] + [B,[C,A]] + [C,[A,B]] = 0$; Lie algebra axiom \\
  & \texttt{casimir\_nonzero}
  & $C_2 \neq 0$ in $\Cl(3,0)$; algebraic spectral gap \\
  & \texttt{casimir\_commutes\_e12}
  & $[C_2, e_{12}] = 0$; Schur centrality \\
\midrule
\texttt{casimir\_eigenvalues.lean}
  & \texttt{casimir\_fund\_val}
  & $C_2 = -2 I_3$ in fundamental of $\so(3)$ \\
  & \texttt{comm\_L1\_L2}
  & $[L_1, L_2] = L_3$; structure constants \\
  & \texttt{casimir\_commutes\_L1}
  & $[C_2, L_1] = 0$; Casimir centrality \\
\midrule
\texttt{casimir\_spectral\_gap.lean}
  & \texttt{killing\_form\_negative}
  & $\Tr(H^2) < 0$ for $H \neq 0$; compactness \\
  & \texttt{so3\_sq\_nonzero}
  & $H \neq 0 \Rightarrow H^2 \neq 0$; spectral gap \\
\midrule
\texttt{block\_tridiagonal.lean}
  & \texttt{grade2\_times\_grade2\_selection}
  & $\pi_1(AB) = \pi_3(AB) = 0$ for bivectors $A, B$ \\
  & \texttt{grade\_decomposition}
  & $x = \pi_0(x) + \pi_1(x) + \pi_2(x) + \pi_3(x)$ \\
\midrule
\texttt{su3\_cartan\_weyl.lean}
  & \texttt{ad\_H1\_eigenvalues\_complete}
  & All 8 eigenvalues of $\ad_{H_1}$ verified \\
  & \texttt{ad\_H2\_eigenvalues\_complete}
  & All 8 eigenvalues of $\ad_{H_2}$ verified \\
  & \texttt{positive\_root\_closed}
  & $\mathfrak{n}^+$ is a closed subalgebra \\
  & \texttt{bracket\_count}
  & All $\binom{8}{2} = 28$ brackets verified \\
  & \texttt{killing\_cartan\_det}
  & $\det B|_{\h} = 108 \neq 0$; semisimplicity \\
\bottomrule
\end{tabular}
\end{table*}

The verification strategy uses explicit finite-dimensional
representations (8-component $\Cl(3,0)$, $3 \times 3$ matrices for
$\so(3)$, 8-component Chevalley basis for $\su(3)$) rather than
mathlib's abstract \texttt{CliffordAlgebra} type. This choice
prioritizes computational transparency: every proof reduces to
\texttt{ext} (extensionality), \texttt{simp} (simplification), and
\texttt{ring} (polynomial identity), making the proofs independently
auditable. A natural extension would connect these concrete types to
mathlib's abstract Clifford algebra via isomorphism proofs.

The full codebase, including all 59 Lean~4 proof files with
over 2,537 verified declarations, is available from the
corresponding author upon reasonable request \cite{uft-github}.

% ============================================================================
\section{Discussion}
\label{sec:discussion}
% ============================================================================

\subsection{The Common Mathematical Parent}

The four results of this paper---two identities and two
analogies---share a common mathematical parent: \emph{character
decomposition on compact groups}. The hierarchy is:

\begin{center}
\begin{tikzcd}[row sep=1em, column sep=0.3em]
  & \text{Peter--Weyl (compact $G$)} \arrow[dl] \arrow[d] \arrow[dr] & \\
  \text{Fortescue ($\ZZ_N$)} &
  \text{Cartan--Weyl ($T^r \subset G$)} &
  \text{Coxeter ($\ZZ_h \subset W$)}
\end{tikzcd}
\end{center}

All three bottom-row decompositions are instances of harmonic analysis
(eigendecomposition under an abelian subgroup) on different subgroups:
\begin{itemize}[leftmargin=*]
\item Fortescue: $\ZZ_N$ (cyclic symmetry of $N$-phase systems).
\item Cartan--Weyl: $T^r$ (maximal torus of a Lie group).
\item Coxeter: $\ZZ_h$ (cyclic group generated by the Coxeter element
  inside the Weyl group $W$).
\end{itemize}

The mass gap / stop band analogy (\Cref{sec:analogy-stopband}) is
connected through the resolvent: the spectral decomposition of an
operator is, at bottom, another instance of decomposing a function
(the resolvent) into contributions from irreducible representations
(eigenspaces).

\subsection{Pedagogical Value}

The Fortescue--Cartan--Weyl connection has immediate pedagogical
applications. Electrical engineering students learn symmetrical
components as a practical tool; the connection to Lie theory reveals
the underlying mathematical structure. Conversely, physics students
learning the Cartan--Weyl decomposition can build intuition from the
power systems parallel: the Cartan subalgebra is the ``zero sequence''
(in-phase, simultaneously measurable quantities), positive roots are
``positive sequences'' (forward-rotating components), and negative
roots are ``negative sequences'' (backward-rotating components).

The correspondence is particularly vivid for $\su(2)$, where the three
generators map directly to the three Fortescue sequences
($N = 3$): $H \leftrightarrow$ zero, $E_+ \leftrightarrow$ positive,
$E_- \leftrightarrow$ negative. The eigenvalues $\{0, +1, -1\}$ match
the three cube roots of unity restricted to
$\{1, \omega_3, \omega_3^2\} \to \{0, +, -\}$.

\subsection{Limitations}

Several limitations should be noted:

\begin{enumerate}[leftmargin=*]
\item \textbf{Concrete vs.\ abstract verification.} Our Lean~4 proofs
  use explicit finite-dimensional representations rather than abstract
  Clifford algebra and Lie algebra APIs. This means the proofs verify
  specific instances ($\so(3)$, $\su(3)$, $\Cl(3,0)$) rather than
  general theorems about all semisimple Lie algebras. The general
  statements (\Cref{prop:peter-weyl-fortescue}, the grade selection rule
  for arbitrary $\Cl(n,0)$) are proved by standard mathematical arguments
  in the text; only the concrete instances are machine-verified.

\item \textbf{The Coxeter analogy.} We have not formalized the Coxeter
  element or principal grading in Lean~4. The Coxeter numbers and
  exponents in \Cref{tab:coxeter} are standard results from
  \cite{humphreys1990}; their machine verification would require a
  formalization of root systems and Weyl groups that is currently beyond
  the scope of our codebase.

\item \textbf{Scope of the grade selection rule.} \Cref{thm:selection-rule}
  is proved for $\Cl(3,0)$. The general statement for $\Cl(n,0)$ follows
  from the standard grading theory of Clifford algebras
  \cite{lounesto2001}, but has not been machine-verified for $n > 3$.

\item \textbf{No physics claims.} We make no claims about quantum field
  theory, the Yang--Mills mass gap, or the physical relevance of
  $\so(14)$. The Casimir eigenvalues and block-tridiagonal structure are
  algebraic facts; their physical interpretation requires additional
  analytical work that is outside the scope of this paper.
\end{enumerate}

\subsection{The Identity--Analogy Boundary}

A contribution of this paper is the explicit delineation between
\emph{identities} and \emph{analogies}:

\begin{table}[htbp]
\centering
\caption{Classification of the four results. Identities are provable
  mathematical equivalences; analogies share structure but have definite
  points of failure.}
\label{tab:identity-analogy}
\begin{tabular}{@{}lll@{}}
\toprule
\textbf{Result} & \textbf{Status} & \textbf{Limitation} \\
\midrule
Peter--Weyl/Fortescue & Identity & None \\
Cartan--Weyl/Fortescue & Identity & None \\
Mass gap/stop band & Analogy & Existence, nonlinearity \\
Coxeter/Fortescue & Analogy & Non-uniform modes \\
\bottomrule
\end{tabular}
\end{table}

The distinction matters. An identity means the two decompositions are
literally the same mathematics, differing only in notation and the
community that uses them. An analogy means the decompositions share a
common mathematical parent (character theory) but differ in scope,
structure, or the level of generality at which they operate.

% ============================================================================
\section{Related Work}
\label{sec:related}
% ============================================================================

The observation that the DFT is the character table of $\ZZ_N$ is
classical (see, e.g., Terras \cite{terras1999}). The connection between
Clifford algebras and Lie algebras is covered in standard references
\cite{lounesto2001,doran2003}. The Coxeter element and principal grading
are treated in Humphreys \cite{humphreys1990} and Kostant
\cite{kostant1959}.

What appears to be new is the systematic identification of Fortescue's
engineering decomposition with the Cartan--Weyl decomposition via their
common parent in character theory, together with the observation that
Clifford grade selection rules impose block-tridiagonal structure on
Hamiltonians. We are not aware of any prior published work that
connects these threads.

Wieser and Song \cite{wieser2021clifford} formalized Clifford algebras
in Lean~3/mathlib, providing abstract algebraic infrastructure; our
proofs use concrete finite-dimensional representations rather than
their abstract type, but a natural extension would connect the two
via isomorphism proofs. Tooby-Smith's PhysLean/HepLean
project \cite{toobysmith2025physlean} formalizes standard physics in
Lean~4 but does not address the cross-domain connections we establish
here.

% ============================================================================
\section{Conclusion}
\label{sec:conclusion}
% ============================================================================

We have established that Fortescue's symmetrical components (1918) and
the Cartan--Weyl root space decomposition (1913--1925) are instances
of the same mathematical operation: character decomposition on compact
groups. The Fortescue matrix is the character table of $\ZZ_N$; the
Cartan--Weyl decomposition is the eigendecomposition of the adjoint
representation under the maximal torus; and the Peter--Weyl theorem
provides the unifying generalization.

We have further shown that Clifford grade selection rules force
operators built from $\so(n)$ generators into block-tridiagonal form,
providing a structural constraint on their spectra. This constraint,
verified for $\Cl(3,0)$ in Lean~4, goes beyond the information encoded
in the Casimir eigenvalues alone.

Two additional analogies---the mass gap as a spectral stop band, and
the Coxeter principal grading as a generalized Fortescue
decomposition---share the same mathematical parent but fall short of
provable identities. We have delineated precisely where each analogy
holds and where it breaks.

All concrete results are machine-verified in Lean~4 with zero
\texttt{sorry} gaps, including the $\so(3)$ structure constants and
Jacobi identity from $\Cl(3,0)$, the Casimir eigenvalue and Killing
form negativity, the grade selection rules, and the complete
Cartan--Weyl decomposition of $\su(3)$ with all 28 independent
brackets.

\begin{acknowledgements}
The machine verification was performed using Lean~4 and the mathlib
library. The author thanks the Lean and mathlib communities for
providing the infrastructure that made this work possible.
Proof engineering and manuscript preparation were assisted by Claude
(Anthropic); all mathematical content, research direction, and
interpretation are the sole responsibility of the author.
\end{acknowledgements}

\section*{Data Availability}
All data generated during this study are contained within the article.
The Lean~4 source code for all machine-verified proofs is available from
the corresponding author upon reasonable request.

% ============================================================================
% References
% ============================================================================

\begin{thebibliography}{30}

\bibitem{fortescue1918}
Fortescue, C.L.:
Method of symmetrical co-ordinates applied to the solution of polyphase
networks.
AIEE Transactions \textbf{37}(2), 1027--1140 (1918)

\bibitem{cartan1913}
Cartan, \'E.:
Les groupes projectifs qui ne laissent invariante aucune multiplicit\'e
plane.
Bulletin de la Soci\'et\'e Math\'ematique de France \textbf{41},
53--96 (1913)

\bibitem{weyl1925}
Weyl, H.:
Theorie der Darstellung kontinuierlicher halbeinfacher Gruppen durch
lineare Transformationen.
Mathematische Zeitschrift \textbf{23}, 271--309 (1925)

\bibitem{peter-weyl1927}
Peter, F., Weyl, H.:
Die Vollst\"andigkeit der primitiven Darstellungen einer geschlossenen
kontinuierlichen Gruppe.
Mathematische Annalen \textbf{97}, 737--755 (1927)

\bibitem{coxeter1934}
Coxeter, H.S.M.:
Discrete groups generated by reflections.
Annals of Mathematics \textbf{35}(3), 588--621 (1934)

\bibitem{kostant1959}
Kostant, B.:
The principal three-dimensional subgroup and the Betti numbers of a
complex simple Lie group.
American Journal of Mathematics \textbf{81}(4), 973--1032 (1959)

\bibitem{humphreys1972}
Humphreys, J.E.:
Introduction to Lie Algebras and Representation Theory.
Graduate Texts in Mathematics, vol.~9. Springer, New York (1972)

\bibitem{humphreys1990}
Humphreys, J.E.:
Reflection Groups and Coxeter Groups.
Cambridge Studies in Advanced Mathematics, vol.~29. Cambridge University
Press (1990)

\bibitem{lounesto2001}
Lounesto, P.:
Clifford Algebras and Spinors, 2nd edn.
London Mathematical Society Lecture Note Series, vol.~286. Cambridge
University Press (2001)

\bibitem{doran2003}
Doran, C., Lasenby, A.:
Geometric Algebra for Physicists.
Cambridge University Press (2003)

\bibitem{hestenes1966}
Hestenes, D.:
Space-Time Algebra.
Gordon and Breach, New York (1966)

\bibitem{georgi1999}
Georgi, H.:
Lie Algebras in Particle Physics, 2nd edn.
Frontiers in Physics. Westview Press (1999)

\bibitem{fulton1991}
Fulton, W., Harris, J.:
Representation Theory: A First Course.
Graduate Texts in Mathematics, vol.~129. Springer, New York (1991)

\bibitem{hall2015}
Hall, B.C.:
Lie Groups, Lie Algebras, and Representations: An Elementary
Introduction, 2nd edn.
Graduate Texts in Mathematics, vol.~222. Springer, Cham (2015)

\bibitem{terras1999}
Terras, A.:
Fourier Analysis on Finite Groups and Applications.
London Mathematical Society Student Texts, vol.~43. Cambridge University
Press (1999)

\bibitem{jaffe-witten2000}
Jaffe, A., Witten, E.:
Quantum Yang--Mills theory.
In: The Millennium Prize Problems, pp.~129--152. Clay Mathematics
Institute (2000)

\bibitem{demoura2021lean4}
de~Moura, L., Ullrich, S.:
The Lean~4 theorem prover and programming language.
In: CADE-28. LNCS, vol.~12699, pp.~625--635. Springer (2021)

\bibitem{mathlib}
The mathlib community:
The Lean mathematical library.
In: CPP 2020, pp.~367--381. ACM (2020)

\bibitem{wieser2021clifford}
Wieser, E., Song, U.:
Formalizing geometric algebra in Lean.
Advances in Applied Clifford Algebras \textbf{31}, 38 (2021)

\bibitem{toobysmith2025physlean}
Tooby-Smith, J.:
HepLean: Digitalising high energy physics.
Computer Physics Communications \textbf{308}, 109457 (2025)

\bibitem{steinmetz1893}
Steinmetz, C.P.:
Complex quantities and their use in electrical engineering.
AIEE Proceedings, 33--74 (1893)

\bibitem{uft-github}
Arnoldy, I.M.:
UFT: Formal verification of unified field theory algebraic scaffold.
\url{https://github.com/iarnoldy/UFT} (2026)

\bibitem{gallier2020}
Gallier, J., Quaintance, J.:
Differential Geometry and Lie Groups: A Computational Perspective.
Geometry and Computing, vol.~12. Springer, Cham (2020)

\bibitem{atiyah-bott-shapiro1964}
Atiyah, M.F., Bott, R., Shapiro, A.:
Clifford modules.
Topology \textbf{3}(Suppl.~1), 3--38 (1964)

\end{thebibliography}

\end{document}
